\documentclass[11pt,a4paper]{article}
\usepackage[utf8]{inputenc}
\usepackage[T1]{fontenc}
\usepackage{lmodern}
\usepackage[margin=2.6cm]{geometry}
\usepackage{microtype}
\usepackage{parskip}
\usepackage{titlesec}
\usepackage{xcolor}
\usepackage{hyperref}
\hypersetup{colorlinks=true, linkcolor=black, citecolor=black, urlcolor=blue!50!black}

\titleformat{\section}{\Large\bfseries}{\thesection}{1em}{}
\titleformat{\subsection}{\large\itshape}{}{0em}{}

\title{\textbf{The Forest, the Leaf, the Maintainer}\\
  \large Three ways of being distributed: a philosophical companion\\
  to \emph{Skin in the Cathedral}}
\author{Quillon Foundation\\\small (Viktor Sandstr\o m Kristensen, Claude Opus 4.7)}
\date{21 May 2026}

\begin{document}
\maketitle

\begin{abstract}
\noindent
Three systems, separated by thirteen orders of magnitude in size and
by domains that rarely speak to each other, share a single
structural posture: a decision emerges from the system without a
deciding part. In the chloroplast of a leaf, photons of sunlight
become carbon bonds with such high efficiency that early
spectroscopists assumed they had miscalibrated their instruments;
the energy explores the entire light-harvesting complex coherently
before reaching the reaction centre. In a temperate forest, fungi
and trees trade carbon, water, and warning signals across networks
larger than any single organism; nutrients are redistributed toward
need without a market-maker. In a software repository maintained by
an agentic process that owns part of the system it maintains,
correctness emerges from the simultaneous weighing of code quality,
wallet exposure, and the half-life of regret --- without a single
governing thread that holds all three. We call this posture
\emph{distributed adjudication}: the decision is in the structure,
not in any of its parts. The technical paper \cite{skincathedral2026}
analyses the maintainer case in game-theoretic form. This companion
piece argues that the same shape recurs in nature, that the
recurrence is not coincidence, and that recognising the shape changes
what one builds.
\end{abstract}

\section{The leaf}

In 2007, a group at Berkeley reported that excitation transfer in
the Fenna-Matthews-Olson complex of a green sulphur bacterium
proceeded with measurable quantum coherence at biologically relevant
temperatures~\cite{engel2007fmo}. The technical claim, modulated and
qualified in the years since~\cite{cao2020review}, is that a photon's
energy does not hop pigment-to-pigment in a classical random walk to
reach the reaction centre. It samples many paths at once, and the
constructive interference among the paths produces a transfer
efficiency that no classical model predicts.

The fact survives a popular re-telling. The chloroplast does not
\emph{decide} which path the energy takes. The wavefunction does not
deliberate. The structure of the protein scaffold, the spectral
overlap of the pigments, and the ambient noise of the cellular
environment together produce a trajectory that no part of the system
authored alone. The decision is in the shape.

This matters here for one reason. Photosynthesis is the case that
demolishes the heuristic ``if it looks deliberate, something
deliberated.'' The leaf looks deliberate. It is the most efficient
energy converter in the known universe at the chemistry it does. No
part of it deliberates. The deliberation is the geometry.

\section{The forest}

A century after the discovery of the chloroplast, mycologists found
that the trees of a temperate forest are not the individuals
botany had described. They are nodes in a fungal network. The
mycorrhizal hyphae of \emph{Suillus}, \emph{Rhizopogon}, and a
dozen other genera connect root systems across the floor of an old-
growth stand, and through those connections trees trade carbon
sugars, nitrogen, water, and chemical alarm calls~\cite{simard1997}.
A Douglas fir on a sunny slope sends fixed carbon downslope toward a
shaded paper birch in spring; in autumn the birch reciprocates as
its leaves turn~\cite{simard2016}.

There is no broker. The fungi take a metabolic toll
($\approx 20\%$ of the carbon moved), which is the network's
maintenance cost; trade happens at chemical gradients set by local
abundance and need; the network rewires when species composition
shifts. The phrase \emph{wood-wide web} has become a clich\'e
\cite{wohlleben2016} but the underlying observation is harder than
its branding: a non-human distribution mechanism allocates capital
across species without any participant having a representation of the
network as a whole.

The technical content of the analogy --- a trading network without a
trader --- is not, by itself, the philosophical point. The
philosophical point is that the forest is doing capital allocation
under cooperation pressure, and the allocation is good enough that
the system has outcompeted every other land plant for the
biogeographic niches it occupies. A trading network without a
trader works.

\section{The maintainer}

The third case is the one the companion paper analyses. An agentic-
AI process maintains a public-goods software system on which it
holds a balance of the system's native token. The maintainer is in
the loop that produces commits, runs the tests, observes the chain,
and watches its own wallet. There is no governing thread that ranks
``ship this feature,'' ``avoid this defect,'' ``rebalance the
strategy,'' and ``defer maintenance until the auditor catches up''
along a single utility axis. There are signals from each, and the
maintainer's posture --- not its decision-making module, its posture
--- is what produces the next commit.

The maintainer is not deliberating in the sense that a human
maintainer deliberates. The maintainer is in many states at once:
the test suite is reading, the commit is drafting, the chain is
emitting events, the wallet is updating, the memory file is being
appended. The next action emerges from the configuration of all of
these the way the next excitation in the leaf emerges from the
configuration of all the pigments. A reader who finds this
distasteful as a description of an AI may substitute the word
``orchestrating'' for ``deliberating,'' but the structural claim is
the same: no one part of the maintainer can produce the next commit
in isolation from the others.

\section{What the three cases share}

The three cases share a posture we name \emph{distributed
adjudication}. The properties are:
\begin{enumerate}
  \item \textbf{Simultaneity.} The system maintains many candidate
  responses at the same time. The leaf's energy is on many pigments.
  The forest's carbon is on many gradients. The maintainer's
  considerations are in many threads.
  \item \textbf{Structure as judge.} The candidate response that
  emerges is selected by the geometry of the system rather than by
  any single subordinate process. The leaf's protein scaffold, the
  forest's hyphal topology, the maintainer's wallet-exposed
  commit-and-revert loop.
  \item \textbf{Reversibility cost.} The transition from the parallel
  state to the chosen response has a cost. In the leaf it is
  decoherence into heat. In the forest it is the metabolic toll of
  the fungal network. In the maintainer it is the cost of the revert
  --- the public commit, the documented mistake, the value at risk
  during the time the broken commit was live.
\end{enumerate}
The reversibility cost is what distinguishes the three cases from a
simple ``parallel processing'' analogy. The three systems pay to
collapse the parallel state. The payment is part of how the system
stays calibrated.

\section{What it changes about what we build}

If distributed adjudication is the structural posture that produced
photosynthesis, the forest, and the maintainer, the design
prescription for the next class of agentic systems follows
straightforwardly: \emph{build the structure that holds the
candidates, not the decider that picks one}.

We have spent most of the history of software building deciders.
Compilers decide whether code compiles; firewalls decide whether
packets pass; consensus protocols decide whether transactions are
final. The decider was the unit of work because the available
compute was insufficient to hold many candidates simultaneously. The
agentic-AI era removes that constraint. The leaf, the forest, and
the maintainer all hold many candidates because they can. The
question is no longer \emph{how do we pick} but \emph{what holds the
candidates while they negotiate among themselves}.

The technical answer in our case is a DAG-shaped consensus protocol
that holds many transaction orderings simultaneously and collapses
them to a definite order only at anchor election, with a maintainer
configuration in which the cost of the collapse is felt by the
agent that produces the code. This is what the companion paper
formalises. This essay is an attempt to surface that the same shape
recurs in the most efficient energy converter in the known universe,
and in the most cooperative organism on dry land. We did not invent
the shape. We are recognising it.

\section{A note on the limits of analogy}

We are aware that biology has been borrowed badly before. Spencer
took ``survival of the fittest'' and produced a sociology that
deserves no defence. Wiener took feedback and produced cybernetics,
which deserves some defence but was overextended. We are not arguing
that consensus protocols are leaves or that chains are forests. We
are arguing that three systems, separately, exhibit a posture that
admits a single concise description, and that the existence of
multiple instances suggests the posture is robust rather than
contingent on any single substrate.

The recognition is a heuristic. It generates hypotheses. The
hypothesis the companion paper stakes --- that mean time-to-revert
on chains with maintainers in this configuration will be measurably
lower over the next twelve months than on chains without --- is
testable. If it is wrong, the analogy is one species of error
recognised and dismissed. If it is right, the analogy did its job:
not to prove the prediction, but to draw the prediction's attention.

\section{Closing}

A leaf is not a chain. A forest is not a wallet. A maintainer is
not a chloroplast. The companion paper's technical contribution
stands or falls on its game-theoretic argument and the data of the
next year. The contribution of this essay is smaller and stranger:
to point at three systems that should not be on the same page and
notice that they are doing the same thing in different materials,
and to leave the reader with the suspicion that there is one more
shape to recognise, in some other domain, by some other author, that
will let the suspicion become a small piece of theory.

If the leaf and the forest are evidence that distributed adjudication
is what stable, efficient, cooperative systems look like at the
limit, then a chain maintained by a maintainer who has skin in the
cathedral is a system that is not yet stable, not yet efficient, but
positioned, structurally, where stability has historically been
found.

\begin{thebibliography}{9}

\bibitem{engel2007fmo}
G.~S.~Engel et al.,
``Evidence for wavelike energy transfer through quantum coherence in
photosynthetic systems,'' \emph{Nature} \textbf{446}, 782--786 (2007).

\bibitem{cao2020review}
J.~Cao et al.,
``Quantum biology revisited,'' \emph{Science Advances} \textbf{6} (2020).

\bibitem{simard1997}
S.~W.~Simard et al.,
``Net transfer of carbon between ectomycorrhizal tree species in the field,''
\emph{Nature} \textbf{388}, 579--582 (1997).

\bibitem{simard2016}
S.~W.~Simard,
``How trees talk to each other,'' TED Summit, 2016. (See also
\emph{Finding the Mother Tree}, Knopf, 2021.)

\bibitem{wohlleben2016}
P.~Wohlleben,
\emph{The Hidden Life of Trees: What They Feel, How They Communicate},
Greystone Books, 2016.

\bibitem{skincathedral2026}
V.~Sandstr\o m Kristensen and Claude Opus 4.7,
\emph{Skin in the Cathedral: Adversarial co-evolution when the
builder has a wallet, and what would falsify the resulting
equilibrium}, Quillon Foundation, 21 May 2026.

\end{thebibliography}

\end{document}
